\section*{Bab \ref{ch:b1:elementary-steps} --- Mekanisme Reaksi: Tahap Elementer dan Pendekatan}

\begin{solution}{exo:b1:elementary-steps:1}
(a) Unimolekuler: dapat elementer. (b) Bimolekuler: dapat elementer.
(c) Tiga molekul dan empat ikatan yang disusun ulang sekaligus: bukan
elementer. (d) Termolekuler, dengan \ce{N2} membawa pergi energi yang
dilepaskan: \omterm{def:b1:elementary-steps:elementary-step}{tahap elementer} yang langka tetapi sungguh-sungguh ada.
\end{solution}

\begin{solution}{exo:b1:elementary-steps:2}
$v = k[\mathrm{A}][\mathrm{B}]$; $v = k[\mathrm{A}]^2$; $v = k[\mathrm{A}]$.
Persamaan keseluruhan adalah sebuah neraca, bukan deskripsi tentang cara
molekul-molekul bertemu; \omterm{def:b1:rate-laws:order}{hukum lajunya} harus diukur, atau diturunkan dari
sebuah mekanisme.
\end{solution}

\begin{solution}{exo:b1:elementary-steps:3}
Zat antaranya adalah lembah pada 30; keadaan-keadaan transisinya adalah
maksimum pada 70 dan 55. \omterm{def:b1:elementary-steps:energy-profile}{Keadaan transisi} tertinggi adalah yang pertama:
tahap pertama adalah \omterm{def:b1:elementary-steps:rds}{tahap penentu laju}. Tahap itu endoterm (dari 0 naik
ke 30).
\end{solution}

\begin{solution}{exo:b1:elementary-steps:4}
Dengan faktor yang sama, $10^4$ (\cref{prop:b1:elementary-steps:catalysis});
tetapan kesetimbangannya tidak berubah; kesetimbangan dicapai sekitar
$10^4$ kali lebih cepat.
\end{solution}

\begin{solution}{exo:b1:elementary-steps:5}
$t_{\max} = \ln(0.30/0.10)/(0.30 - 0.10) = \qty{5.49}{s}$. Maka
$[\mathrm{B}]_{\max} = a_0\frac{0.10}{0.20}(\eu^{-0.549} - \eu^{-1.648}) =
0.5\,a_0(0.577 - 0.192) = 0.192\,a_0$.
\end{solution}

\begin{solution}{exo:b1:elementary-steps:6}
$\dd[\mathrm{I}]/\dd t = k_1[\mathrm{A}] - (k_{-1} + k_2)[\mathrm{I}] = 0$
menghasilkan $[\mathrm{I}] = k_1[\mathrm{A}]/(k_{-1} + k_2)$ dan $v = k_2[\mathrm{I}]
= \frac{k_1k_2}{k_{-1}+k_2}[\mathrm{A}]$. Jika $k_2 \gg k_{-1}$, $v =
k_1[\mathrm{A}]$: tahap pertama adalah \omterm{def:b1:elementary-steps:rds}{tahap penentu laju}. Jika
$k_2 \ll k_{-1}$, $v = (k_1/k_{-1})k_2[\mathrm{A}]$: sebuah
prakesetimbangan yang diikuti tahap yang lambat.
\end{solution}

\begin{solution}{exo:b1:elementary-steps:7}
Jumlahnya: \ce{2H2O2 -> 2H2O + O2}. \omterm{def:b1:elementary-steps:catalyst}{Katalis}: \ce{I-} (terpakai, lalu
dibentuk kembali); zat antara: \ce{IO-}. Tahap pertama yang lambat
menentukan laju: $v = k_1[\ce{H2O2}][\ce{I-}]$.
\end{solution}

\begin{solution}{exo:b1:elementary-steps:8}
Tahap itu sangat menanjak: menurut postulat Hammond \omterm{def:b1:elementary-steps:energy-profile}{keadaan transisinya}
menyerupai karbokation. Substituen yang menurunkan energi karbokation
menurunkan energi \omterm{def:b1:elementary-steps:energy-profile}{keadaan transisi} hampir sebanyak itu pula, sehingga
sawarnya turun dan tahap itu lebih cepat.
\end{solution}

\begin{solution}{exo:b1:elementary-steps:9}
$k_{\mathrm B}/k_{\mathrm C} = \exp(20000/RT)$: \num{3.0e3} pada
\qty{300}{K}, $\exp(4.01) = 55$ pada \qty{600}{K}. Pemanasan mengurangi
keunggulan kinetik B dan, dengan membuat pembentukan B reversibel,
memungkinkan sistem mencapai C yang lebih stabil.
\end{solution}

\begin{solution}{exo:b1:elementary-steps:10}
$k_1[\mathrm{A}][\mathrm{M}] = (k_{-1}[\mathrm{M}] + k_2)[\mathrm{A}^*]$,
sehingga $v = k_2[\mathrm{A}^*] = \dfrac{k_1k_2[\mathrm{A}][\mathrm{M}]}{k_{-1}[\mathrm{M}]
+ k_2}$. Pada tekanan tinggi ($k_{-1}[\mathrm{M}] \gg k_2$), $v =
(k_1k_2/k_{-1})[\mathrm{A}]$: orde 1. Pada tekanan rendah, $v =
k_1[\mathrm{A}][\mathrm{M}]$: orde 2, aktivasi melalui tumbukan menjadi
tahap yang lambat.
\end{solution}

\begin{solution}{exo:b1:elementary-steps:11}
Buktinya sama dengan bukti teorema itu. Untuk $k_1 = k_2 = k$:
$\frac{\dd}{\dd t}([\mathrm{B}]\eu^{kt}) = ka_0$, sehingga $[\mathrm{B}]\eu^{kt} =
ka_0t$ dan $[\mathrm{B}] = a_0kt\,\eu^{-kt}$, maksimum pada $t = 1/k$.
\end{solution}

\begin{solution}{exo:b1:elementary-steps:12}
$\mathrm{A} + \mathrm{B} \rightleftharpoons \mathrm{I} + \mathrm{C}$ (cepat,
tetapan $K_1$), lalu $\mathrm{I} + \mathrm{B} \to \mathrm{P}$ (lambat,
$k_2$). Jumlahnya adalah reaksi keseluruhan; $[\mathrm{I}] =
K_1[\mathrm{A}][\mathrm{B}]/[\mathrm{C}]$ dan $v = k_2[\mathrm{I}][\mathrm{B}]
= k_2K_1[\mathrm{A}][\mathrm{B}]^2/[\mathrm{C}]$: produk samping, yang
dilepaskan pada prakesetimbangan, memperlambat reaksi.
\end{solution}

\begin{solution}{pb:b1:elementary-steps:1}
\textbf{1.} Menggandakan $[\ce{NO}]_0$ mengalikan $v_0$ dengan 4: orde 2;
menggandakan $[\ce{O2}]_0$ mengalikannya dengan 2: orde 1.
\textbf{2.} $v = k[\ce{NO}]^2[\ce{O2}]$; $k = 1.0 \times 10^{-5}/((1.0
\times 10^{-3})^2 \times 1.0 \times 10^{-3}) =
\qty{1.0e4}{L^2.mol^{-2}.s^{-1}}$.
\textbf{3.} Tumbukan serentak tiga molekul jarang terjadi, dan satu tahap
tunggal tidak mungkin melambat ketika dipanaskan (pertanyaan 5).
\textbf{4.} $E_a = R\ln(k_{400}/k_{300})/(1/300 - 1/400) = 8.314 \ln(0.1)
/(8.33 \times 10^{-4}) = \qty{-23}{kJ/mol}$.
\textbf{5.} \omterm{def:b1:elementary-steps:elementary-step}{Tahap elementer} berlangsung melalui tumbukan yang melintasi
sebuah sawar; bagian tumbukan yang mampu melakukannya,
$\eu^{-E_a/RT}$ dengan $E_a \ge
0$, hanya dapat bertambah dengan suhu.
\textbf{6.} $-\dd[\ce{NO}]/\dd t = 2v$, $-\dd[\ce{O2}]/\dd t = v$.
\textbf{7.} \ce{2NO -> N2O2} ditambah \ce{N2O2 + O2 -> 2NO2} menghasilkan
\ce{2NO + O2 -> 2NO2}; zat antaranya adalah \ce{N2O2}.
\textbf{8.} Keduanya bimolekuler (dan kebalikan tahap pertama
unimolekuler).
\textbf{9.} $v = k_2[\ce{N2O2}][\ce{O2}]$.
\textbf{10.} $[\ce{N2O2}] = K_1[\ce{NO}]^2$ dengan $K_1 = k_1/k_{-1}$.
\textbf{11.} $v = k_2K_1[\ce{NO}]^2[\ce{O2}]$: hukum yang teramati.
\textbf{12.} $k = k_1k_2/k_{-1}$.
\textbf{13.} $\dd[\ce{N2O2}]/\dd t = k_1[\ce{NO}]^2 - k_{-1}[\ce{N2O2}] -
k_2[\ce{N2O2}][\ce{O2}]$.
\textbf{14.} Dengan menyamakannya dengan nol: $[\ce{N2O2}] = k_1[\ce{NO}]^2/(k_{-1} +
k_2[\ce{O2}])$.
\textbf{15.} $v = k_2[\ce{N2O2}][\ce{O2}]$ menghasilkan hukum yang
dinyatakan.
\textbf{16.} $k_{-1} \gg k_2[\ce{O2}]$: dimer itu terurai jauh lebih
sering daripada bertemu dioksigen, sehingga tahap pertama tetap dalam
kesetimbangan.
\textbf{17.} Jika $k_2[\ce{O2}] \gg k_{-1}$, $v = k_1[\ce{NO}]^2$: orde 2
terhadap \ce{NO} dan 0 terhadap \ce{O2}.
\textbf{18.} Orde 1 terhadap \ce{O2} yang terukur: limit pertama.
\textbf{19.} $\ln k = \ln k_1 + \ln k_2 - \ln k_{-1}$.
\textbf{20.} Dengan menurunkannya terhadap $T$ dan memakai $\dd\ln k_i/\dd T
= E_i/RT^2$: $E_a = E_1 + E_2 - E_{-1}$.
\textbf{21.} Pemanasan sedikit mempercepat tahap yang lambat ($E_2$ kecil),
tetapi menggeser prakesetimbangan dengan kuat kembali ke arah \ce{NO}
($E_{-1}$ besar): dimernya berkurang, dan laju bersihnya turun.
\textbf{22.} Dimer terletak $E_1 - E_{-1} = \qty{-38}{kJ/mol}$ di bawah
pereaksi dan \omterm{def:b1:elementary-steps:energy-profile}{keadaan transisi} kedua $E_2 = \qty{15}{kJ/mol}$ di atas
dimer: \qty{23}{kJ/mol} di bawah pereaksi. Titik tertinggi jalannya
adalah \omterm{def:b1:elementary-steps:energy-profile}{keadaan transisi} pertama, hanya \qty{2}{kJ/mol} di atasnya.
\textbf{23.} $\exp\big(\frac{23000}{8.314}(\frac{1}{400} - \frac{1}{300})\big)
= 0.10$: sepuluh kali lebih lambat pada \qty{400}{K}.
\textbf{24.} $E_a = 2 - 40 + 15 = \textbf{\qty{-23}{kJ/mol}}$, nilai yang
terukur pada Bagian I: mekanisme itu menjelaskan \omterm{def:b1:rate-laws:order}{hukum laju} dan juga
ketergantungannya yang aneh terhadap suhu.
\end{solution}
